
\documentclass[12pt,  % 12 Punkt Schrift
  a4paper,            % DIN A4
  oneside,            % einseitiger Druck
  BCOR=5mm,           % Bindungskorrektur
  parskip=half,       % Halber Abstand zwischen Absätzen
]{scrreprt}
\usepackage{./style/thesisstyle}
\usepackage{enumitem}           % customizations for enumerations
\usepackage{listings}   				% für Codeschnipsel in lstlisting{}
\usepackage{siunitx}

\definecolor{newgray}{rgb}{0.92, 0.92, 0.92}

\begin{document}
% !TEX root = ../VPJ.tex
%
% configurations
%

% text field
%-> replace supervisor names with correct ones
\firstSupervisor{Prof. Dr.-Ing. Ulfert Meiners}
\secondSupervisor{Prof. Dr.-Ing. Jochen Maaß, Prof. Dr.-Ing. Michael Röther}

% text field
%-> replace title with your thesis title
\thesisTitle{Verbundprojekt 2018/2019}
\thesisTitleEN{Verbundprojekt 2018/2019}

%\thesisKind{VPJ}

% text field
%-> replace the key words with your own key words
\keywordsDE{Smart Kamera, Objekterkennung, Farbdetektion, Gehäusekonstruktion, PixyCam, Raspberry Pi}
\keywordsEN{Smart Camera, Objecttracking, Colordetection, Construction, PixyCam, Raspberry Pi}

% text field
%-> replace the text with a description of the thesis
\abstractDE{
Im Rahmen dieser Arbeit werden zwei Kamerasysteme verglichen, welche farbige Objekte auf einem Laufband erkennen und deren Farbe auswerten. Über verschiedene Schnittstellen wird die Auswertung ausgegeben. Dazu werden Anwendungen entworfen, welche die Auswertung auf einem der Systeme ermöglichen. Die Arbeit beschreibt außerdem die konkrete Implementierung des Quellcodes der entworfenen Anwendungen für eines der Kamerasysteme. Ein weiterer zentraler Punkt der Arbeit ist der Entwurf einer geeigneten Halterung für die Kamerasysteme.} 
\abstractEN{
In this thesis, two camera systems are compared, which recognize colored objects on a conveyor belt and evaluate their color. The evaluation is output via various interfaces. For this purpose, applications are designed, which enable the evaluation on one of the systems. The thesis also describes the implementation of the source code of the designed applications for one of the camera systems. Another point of the work is the design of a suitable holder for the camera systems.
} 

% text field
%-> replace jon with your name
\thesisAuthor{Daniel Friedrich}

% text field
%-> enter the submission date
\submissionDate{25. Mai 2019}

% switch - uncomment only one
%-> uncomment NDA or public
%\NDA{yes}
\NDA{no}

% switch - uncomment only one
%-> uncomment the study course you are in
%\studycourse{TI}
%\studycourse{AI}
%\studycourse{WI}
%\studycourse{BMT}
%\studycourse{MAI}
\studycourse{MAA}
    % load all settings
\newcommand{\inlinetodo}[1]{\todo[inline,linecolor=blue,backgroundcolor=blue!25,bordercolor=blue,size=\normalsize]{#1}}

\hyphenation{Ba-che-lor-the-sis Mas-ter-the-sis Auf-trags-List-Wid-get item-Double-Clicked}

% Cover page here, no page number
% !TEX root = ../VPJ.tex
%
% cover page
% @author Thomas Lehmann
%

\thispagestyle{empty}
\begin{titlepage}
{\fontfamily{phv}\selectfont

% NDA, if needed
% TODO: Select correct font
\begin{textblock*}{\textwidth}(75mm,9mm)
  \begin{minipage}[b][0cm][b]{\textwidth}
  \fontsize{16pt}{16pt}
  \selectfont
    \begin{flushleft}
      \IthesisNDAFull
    \end{flushleft}
  \end{minipage}
\end{textblock*}

% kind of thesis
% TODO: Select correct font
\begin{textblock*}{\textwidth}(30mm,115mm)
  \begin{minipage}[b][0cm][b]{\textwidth}
    \fontsize{22pt}{20pt}
    \selectfont
  	\begin{flushright}
      \IthesisKind
  	\end{flushright}
  \end{minipage}
\end{textblock*}

% author of thesis
% TODO: Select correct font
\begin{textblock*}{\textwidth}(30mm,140mm)
  \begin{minipage}[b][0cm][b]{\textwidth}
  \fontsize{14pt}{20pt}
  \selectfont
    \begin{flushright}
      \IthesisAuthor \\
			\IthesisAuthorSecond
  	\end{flushright}
  \end{minipage}
\end{textblock*}

% Title of thesis
% TODO: Select correct font
\begin{textblock*}{\textwidth}(30mm,155mm)
  \begin{minipage}[b][0cm][b]{\textwidth}
  \fontsize{18pt}{20pt}
  \selectfont
  	\begin{flushright}
       \IthesisTitle
  	\end{flushright}
  \end{minipage}
\end{textblock*}

% german version of faculty and department
\begin{textblock*}{\textwidth}(20mm,256mm)
  \begin{minipage}[b][0cm][t]{\textwidth}
  \fontsize{11pt}{10pt}
  \selectfont
    \begin{flushleft}
      \textit{\IthesisFacultyFull} \\
      \textit{\IthesisDepartmentFull}
    \end{flushleft}
  \end{minipage}
\end{textblock*}

% english version of faculty and department
\begin{textblock*}{\textwidth}(52mm,256mm)
  \begin{minipage}[b][0cm][t]{\textwidth}
  \fontsize{11pt}{10pt}
  \selectfont
    \begin{flushright}
      \textit{\IthesisFacultyFullEN} \\
      \textit{\IthesisDepartmentFullEN}
    \end{flushright}
  \end{minipage}
\end{textblock*}
}
\end{titlepage}
\                % this backslash is needed, otherwise LaTeX does wired things ....


% Titlepage is page one even if the number is not shown.
\pagenumbering{roman}
% Title page here
% !TEX root = ../VPJ.tex
%
% title page
% @author Thomas Lehmann
% Hints for titel page and page numbering: https://en.wikipedia.org/wiki/Title_page
%
\newpage
\thispagestyle{empty}
{\fontfamily{phv}\selectfont

  % Author of thesis
  % TODO: Select correct font
  \vspace*{1cm}
  \begin{minipage}[b]{\textwidth}
    \fontsize{14pt}{20pt}
    \selectfont
    \begin{center}
      \IthesisAuthor
    \end{center}
  \end{minipage}

  % Title of thesis
  % TODO: Select correct font  
  \vspace{1.5cm}
  \begin{minipage}[b][0cm][b]{\textwidth}
    \fontsize{18pt}{20pt}
    \selectfont
    \begin{center}
      \IthesisTitle  	  
  	\end{center}
  \end{minipage}

  % Important information
  % TODO: Select correct font  
  \begin{textblock*}{\textwidth}(40mm,210mm)
    \begin{minipage}[b]{\textwidth}
      \fontfamily{ptm}\selectfont
      \fontsize{12pt}{14pt}
      \selectfont
      \IthesisKindDE ~eingereicht im Rahmen der \IthesisExaminationDE \\

      im Studiengang \IstudyCourseName \\
      am \IthesisDepartmentFull \\
      der Fakultät Technik und Informatik\\
      der Hochschule für Angewandte Wissenschaften Hamburg\\

      Betreuender Prüfer: \IfirstSv \\

      Eingereicht am: \ISubDate \\
    \end{minipage}
  \end{textblock*}
}


% Abstract page here
% !TEX root = ../VPJ.tex
%
% abstract page
% @author Thomas Lehmann
%
\newpage
\thispagestyle{empty}   % TODO: This page should show page number
\clearpage
\textbf{\IthesisAuthor \hspace{0.5cm}  \IthesisAuthorSecond}

\vspace{0.5cm}
\textbf{Thema der Arbeit}

\IthesisTitle

\vspace{0.3cm}
\textbf{Stichworte}

\IkeyWordsDE

\vspace{0.3cm}
\textbf{Kurzzusammenfassung}

\begin{minipage}{\textwidth}
\IabstractDE
\end{minipage}

\vspace{1.0cm}
\textbf{\IthesisAuthor \hspace{0.5cm} \IthesisAuthorSecond}

\vspace{0.3cm}
\textbf{Title of Thesis}

\IthesisTitleEN

\vspace{0.3cm}
\textbf{Keywords}

\begin{minipage}{\textwidth}
\IkeyWordsEN
\end{minipage}

\vspace{0.3cm}
\textbf{Abstract}

\IabstractEN


\newpage

  \vspace{1.5cm}
  \begin{minipage}[b][0cm][b]{\textwidth}
    \fontsize{22pt}{20pt}
		\selectfont
    \begin{center}
      \textbf{Glossar}  	  
  	\end{center}
  \end{minipage}

\begin{table}[!ht]
	\centering
	\renewcommand{\arraystretch}{1.5}
	\begin{tabular}{|c|p{10cm}|}
		\hline
		\textbf{Begriff} &	\textbf{Beschreibung} \\
		\hline
    Auftrag & Ein Auftrag ist das vom Benutzer eingegebene Konstrukt aus verschiedenen Prozessen \\
		\hline
    Arbeitsplatz & Ein einzelner Bearbeitungsplatz an einer Station\\
		\hline
    Einhorn & Eigenname von Robotino 5  \\
		\hline
    Fertigungsstraße & Zusammenfassung aller Stationen und Arbeitsplätze  \\
		\hline
    Fertigungsrechner & Rechner mit Datenbank und Kommunikation zu RFID  \\
		\hline
    Fertigungsplanungsrechner & Rechner mit UDP-Kommunikation und Auftragsplanung  \\
		\hline
    Position & Enthält eine Kodierung einer Station oder Arbeitsplatzes \\
		\hline
		Produktionsprozess & Ein Produktionsprozess ist im Zusammenhang mit der Datenbank als Rezept für die Bearbeitung des Werkstücks zu verstehen. \\
		\hline
    Prozess & Ein Prozess beschreibt die Abfolge unterschiedlich vieler Prozessschritte. Im Zusammenhang der Datenbank wird der Begriff Werkstück verwendet  \\
		\hline
    Prozessschritt & Enthält genau einen Arbeitsplatz mit Bearbeitungsdauer \\
		\hline
    Station & Der Zusammenschluss aus einem RFID Lese- Schreibgerät und den dahinterliegenden beiden Arbeitsplätzen \\
		\hline
    Werkstück & Ein Werkstück symbolisiert exakt einen "`Stein"' in der Fertigungsstraße. Ein Werkstück entspricht immer genau einem Prozess\\
		\hline
	\end{tabular}
	\label{tab:Glossar}
\end{table}

% Table of contents here
\tableofcontents


% path to the chapters folder is set to find the images used there
\graphicspath{ {./chapters/} }

\setcitestyle{square}

% Chapters
\clearpage
\pagenumbering{arabic}
% !TEX root = ../VPJ.tex

\chapter{Einleitung}
\label{sec:Einleitung}

Um an einem realen System Abläufe und Prozesse zu simulieren, wird den Studierenden eine Anlage zur Verfügung gestellt, welche hauptsächlich aus Teilen der Firma „Festo“ besteht (folgend: Festo-Transfersystem). An diesem Festo-Transfersystem können die Studierende Laufbänder, Ampelanlage, Lichtschranken, Weichen etc. mithilfe selbst geschriebener Programme ansteuern und auslesen. 

Das Laufband kann aufgelegte Objekte transportieren. Dabei kann die aktuelle Position des Objekts durch diverse Lichtschranken erfasst werden. Am Anfang des Laufbands kann mithilfe eines Entfernungssensors die Höhe des Objekts erfasst werden. Mit den vorhandenen Komponenten kann das Festo-Transfersystem allerdings ausschließlich die Position und die Höhe der Objekte bestimmen und nicht deren Farbe, Form, Ausrichtung oder sonstige optische Merkmale. 

\begin{figure}[htb]
    \centering
    \includegraphics[width=0.6\textwidth]{Abbildungen/Platzhalter.jpg}
    \caption[CMUcam5 Pixy \newline Quelle: \textit{http://charmedlabs.com/default/products/}]{CMUcam5 Pixy}		
    \label{fig:PixyCam}
\end{figure}

Eine Erweiterung um ein Kamerasystem kann diese Mängel beheben und sogar noch weitere Funktionen und Features hinzufügen. 

\section{Vorhaben}
\inlinetodo {blabla}
Ähnlich einer Smart Kamera \cite{festo_dok}, bei der nicht nur das Kamerabild, sondern auch weitere schon verarbeitete Informationen ausgegeben werden sollen, wird ein Prozessor benötigt, welcher in der Lage ist Bildverarbeitung durchzuführen. Es wurde als Möglichkeit für ein Prozessorsystem empfohlen, sich mit dem BeagleBone Black und Raspberry PI auseinanderzusetzen. Mithilfe einer Marktrecherche werden zunächst weitere geeignete Kamerasysteme gesucht und anschließend anhand der Eignung sortiert. Dabei wird auch eine geeignete Schnittstelle zu dem Kamerasystem ausgewählt, mit dem die Bilder und Informationen an den PC gesendet werden. Nachdem die beiden geeignetsten Kamerasysteme gewählt und bestellt sind, werden zugehörige Gehäuse und Befestigungen konstruiert und gefertigt. Nebenbei werden erste einfache Testprogramme geschrieben, um die Funktion von den gewählten Kamerasystemen zu gewährleisten. Der erste Prototyp wird dabei voraussichtlich mit Hilfe rapid-prototyping ein Erzeugnis aus dem 3D-Druck sein, um Kamera an der richtigen Position zu halten. Nachdem das Endprodukt montiert ist, folgt eine Auswertung der Kamerabilder. Folgend folgen optional das Erstellen einer Bibliotheksdatei und eine Automation der Kalibrierung und des Weißabgleichs. Der Zugriff auf die Kamera soll gewährleistet sein. 

\begin{figure}[htb]
    \centering
    \includegraphics[width=0.4\textwidth]{Abbildungen/Platzhalter.jpg}
    \includegraphics[width=0.4\textwidth]{Abbildungen/Platzhalter.jpg}
    \caption{PixyCam Gehäusefehler}		
    \label{fig:Pixy_Fehler}
\end{figure}


% !TEX root = ../VPJ.tex

\chapter{Grundlagen und Konzepte}
\label{sec:Grundlagen}

In diesem Kapitel werden einige Grundlagen erläutert, die zum Verständnis der späteren Kapitel der Arbeit beitragen. 

\section{Qt}

Qt ist eine plattformübergreifende Anwendungsentwicklungsumgebung für Desktopanwendungen, eingebettete und mobile Systeme. Unterstützt werden unter anderem Linux, OS X,  Windows, Android, iOS, BlackBerry oder Sailfish OS. 

Die Umgebung ist in C++ geschrieben, weshalb vorwiegend in dieser Sprache programmiert wird. Eine Besonderheit an Qt ist die Benutzung eines Signal-Slot-Systems. 

Ähnlich anderen Entwicklungsumgebungen bietet Qt mit einem implementierten Designer Möglichkeiten zur GUI-Programmierung mit vorgefertigten Widgets (vgl. \cite{qt_allgemein}).

Die nächsten Abschnitte zeigen einige Besonderheiten von Qt und Konzepte, die in der Ausarbeitung des Projekts verwendet wurden. 

\subsection{Eventbasiertes System}
\label{sec:Eventbasiert}

Innerhalb von Qt werden Signale und Slots zur Kommunikation zwischen Objekten genutzt. Das Signal-Slot Konzept ersetzt das in anderen Umgebungen verwendete Callback-System, welches beispielsweise in Eclipse bei der Java-Programmierung oder in Visual Studio mit C\# verwendet wird. 

\begin{figure}[htb]
    \centering
    \includegraphics[width=0.5\textwidth]{Abbildungen/signalslots.png}
    \caption[Signale und Slots in Qt \newline Quelle: \textit{https://doc.qt.io/qt-5/signalsandslots.html}]{Signale und Slots in Qt}		
    \label{fig:signalslots}
\end{figure}

In Abbildung \ref{fig:signalslots} ist das Konzept visuell dargestellt. Ein Signal wird immer dann emittiert, wenn ein bestimmtes Event auftritt. Die in Qt verwendeten Widgets stellen bereits viele dieser Signale bereit, können jedoch um eigene ergänzt werden. 

Das Gegenstück zu den Signalen sind Slots, welche als Antwort auf ein Signal aufgerufen wird. Es können beliebig viele Signale an einen Slot gebunden werden. Es ist ebenfalls möglich Signale direkt an weitere Signale zu binden. 

Ein Slot kann wie eine normale Funktion auch mittels Funktionsaufruf verwendet werden.  

Nähere Informationen zu diesem Konzept, zu Signalen und zu Slots können der Quelle \cite{qt_signalslot} entnommen werden. 

Um Signale und Slots zu verbinden werden verschiedene Möglichkeiten bereitgestellt. In dieser Arbeit wurde hauptsächlich mit der connect() Funktion gearbeitet, die von jeder Klasse, die von QObject erbt, bereitgestellt wird. 

\begin{lstlisting}[frame=single, breaklines=true, numbers=left, stepnumber=2, firstnumber=1, numberstyle = \tiny, caption=Connect Signals and Slots ,label=lst:Connect]
QObject::connect(const QObject *sender, const char *signal, const QObject *receiver, const char *method, Qt::ConnectionType type = Qt::AutoConnection)
\end{lstlisting}

In Listing \ref{lst:Connect} ist der Connect-Aufruf vollständig dargestellt. Hinzuzufügen ist, dass der letzte Übergabeparameter, ConnectionType, optional ist, da er als Standardwert QT::AutoConnection beinhaltet. Die ersten beiden Parameter enthalten also immer die Senderseite, die letzten beiden Parameter die Empfängerseite. Dabei ist zuerst das Objekt anzugeben, welches sendet oder empfängt und danach das Signal, bzw. der Slot, der genutzt werden soll. 

\subsection{Sockets}
\label{sec:QTSocket}

Die Kommunikation über UDP wird in Qt mit einem QUdpSocket bereitgestellt. Der Socket bietet Funktionen, mit denen UDP-Datagramme gesendet und empfangen werden können. 

Um einen Datenaustausch zu ermöglichen wird der Socket mittels einer Bind()-Methode an eine spezielle Adresse und einen Port gebunden. Vorteil dieser Implementierung ist, dass beim Erhalt einer Nachricht ein readyRead()-Signal emittiert wird, auf das sich verbunden werden kann. 

Sobald das readyRead() Signal emittiert wurde liefert hasPendingDatagrams() true zurück. Das Datagramm kann anschließend über readDatagram() oder receiveDatagram() ausgelesen werden (vgl. \cite{qt_socket}).

Außerdem unterstützt die Klasse UDP multicast.

Wenn ein Datagram nicht ausgelesen wird, wird kein weiteres readyRead-Signal emittiert, bis dieses ausgelesen wurde. 

\subsection{Datenbank in Qt}
\label{sec:QTDatabase}

Um möglichst Einfach Datenbanken in Qt zu implementieren wird die Klasse QSqlDatabase (vgl. \cite{qt_database}) bereitgestellt. Diese Klasse ist für die Verbindung und Kommunikation an eine bestehende Datenbank zuständig. Der Datenbanktyp wird über den eingeladenen Treiber bestimmt. In dieser Arbeit sollte auf eine Datenbank vom Typ MYSQL zugegriffen werden.

Um den Treiber der Entwicklungsumgebung von Qt bekannt zu machen war es notwendig eine MYSQL Library einzubinden. Dies konnte gelöst werden, indem die Datei "`MySQLLib.dll"', welche zum Beispiel im MySQL Connector zu finden ist (Pfad: C:\textbackslash Program Files\textbackslash MySQL\textbackslash MySQL Connector.C 6.1\textbackslash lib), an einen Ort kopiert wird, dessen Pfad von der Entwicklungsumgebung erkannt wird. Ein Möglicher Ort ist C:\textbackslash Windows. 

Nachdem sich mit dem Setzten von Nutzername, Passwort, Datenbankname und Host auf die Datenbank verbunden wurde können die gegebenen Funktionen verwendet werden. 

In dieser Arbeit wurde eine QSqlQuery innerhalb der QSqlDatabase verwendet, um die Daten der Datenbank auszulesen und zu manipulieren. Die Query stellt Funktionen zur Verfügung, mit denen die gängigen SELECT, INSERT und UPDATE Befehle umgesetzt werden können.  

\subsection{State Machines in Qt}
\label{sec:StateMachines}

Um entwickelte Zustandsautomaten zu implementieren bietet Qt ein Konzept "`State Machines"' mit dazugehörigen Klassen. 

Dieses Konzept basiert darauf, einen in beispielsweise UML vorliegenden Zustandsgrafen möglichst einfach zu implementieren. 

Eine einfache Zustandsmaschine, dargestellt in Abbildung \ref{fig:StateMachineexample}, kann wie folgt implementiert werden:

\begin{figure}[htb]
    \centering
    \includegraphics[width=0.7\textwidth]{Abbildungen/statemachinebutton.png}
    \caption[State Machine Beispiel \newline Quelle: \textit{https://doc.qt.io/qt-5/statemachine-api.html}]{State Machine Beispiel}		
    \label{fig:StateMachineexample}
\end{figure}

Zunächst wird die State Machine und ihre Zustände erzeugt (Listing \ref{lst:SM1}). 

\begin{lstlisting}[frame=single, breaklines=true, numbers=left, stepnumber=2, firstnumber=1, numberstyle = \tiny, caption=State Machine Beispiel Teil 1 ,label=lst:SM1]
    QStateMachine machine;
    QState *s1 = new QState();
    QState *s2 = new QState();
    QState *s3 = new QState();
\end{lstlisting}

Anschließend werden jedem Zustand alle für diesen Zustand vorhandenen Transitionen mit der QState::addTransition() Funktion hinzugefügt (Listing \ref{lst:SM2}). Dabei wird angegeben, bei welchem Event (Signal) die Transition ausgeführt werden soll. Wenn ein Transitionssignal auserhalb des Zustands auftritt wird die Transition nicht ausgeführt. 

\begin{lstlisting}[frame=single, breaklines=true, numbers=left, stepnumber=2, firstnumber=1, numberstyle = \tiny, caption=State Machine Beispiel Teil 2 ,label=lst:SM2]
    s1->addTransition(button, SIGNAL(clicked()), s2);
    s2->addTransition(button, SIGNAL(clicked()), s3);
    s3->addTransition(button, SIGNAL(clicked()), s1);
\end{lstlisting}

Listing \ref{lst:SM3} zeigt, wie anschließend die Zustände dem Zustandsautomaten hinzugefügt werden. Außerdem wird der initiale Zustand gesetzt, der bei Programmstart aktiv sein soll. 

\begin{lstlisting}[frame=single, breaklines=true, numbers=left, stepnumber=2, firstnumber=1, numberstyle = \tiny, caption=State Machine Beispiel Teil 3 ,label=lst:SM3]
    machine.addState(s1);
    machine.addState(s2);
    machine.addState(s3);
    machine.setInitialState(s1);
\end{lstlisting}

Abschließend muss die Zustandsmaschine gestartet werden. Ab diesem Zeitpunkt reagiert der Automat auf die Transitionen und wechselt zwischen den Zuständen. 

\begin{lstlisting}[frame=single, breaklines=true, numbers=left, stepnumber=2, firstnumber=1, numberstyle = \tiny, caption=State Machine Beispiel Teil 4 ,label=lst:SM4]
    machine.start();
\end{lstlisting}

Hauptvorteil der so implementierten Zustandsmaschine ist, dass diese vollkommen aysnchron zum restlichen Programm abläuft- somit müssen keine zusätzlichen Schleifen in Tasks gestartet werden. 

Um Funktionen mit einem Zustand zu verknüpfen kann eine Signal, welches jeder Zustand emittiert genutzt werden. Beim Eingang in einen Zustand wird das Signal entered() emittiert, beim verlassen eines Zustands das Signal exited(). In dieser Arbeit wurde ausschließlich mit dem entered() Signal gearbeitet, da dies die Anforderungen ausreichend erfüllt hat. 

Das aufgezeigte Beispiel und weitergehende Erklärungen zu State Machines und den weiteren Klassen sind in \cite{qt_statemachine} zu finden. Darunter sind auch Konzepte zu finden, die ermöglichen Eigenschaften direkt an einen Zustand zu binden. 

\subsection{QListWidget}
\label{sec:QListWidgetItem}

Das QListWidget ermöglicht es, selbst programmierte QListWidgetItems in einer Liste zu organisieren und darzustellen. Zusätzlich werden diverse Funktionen bereitgestellt, mit denen sich die QListWidgetItems ordnen, ergänzen oder finden lassen (vgl. \cite{qt_listwidget}). 

Über die Auswahl der allgemeinen Widgets kann das QListWidget im Layout-Manager ausgewählt werden. 

% !TEX root = ../VPJ.tex

\chapter{Schnittstellen}
\label{sec:Schnittstellen}

\inlinetodo{Schnittstellen}

\section{Kommunikation zu Robotinos}
\label{sec:Gewerk2Protokoll}

Die Robotinos werden von Gewerk 2 verwaltet und programmiert. Um eine erfolgreiche Kommunikation zwischen Fertigungsplanungsrechner und Robotino zu gewährleisten existieren folgende Vereinbarungen mit Gewerk 2.

\subsection{Sequenzdiagramm}
\label{sec:sequenzdiagram}

Zunächst wurde eine Struktur entwickelt, in der beschrieben wird, wie die Daten zwischen Robotino und Fertigungsplanungsrechner ausgetauscht werden sollen. Dies wurde anschließend in einem Sequenzdiagramm, Abbildung \ref{fig:Sequenzdiagramm}, dargestellt. 

\begin{figure}[htb]
    \centering
    \includegraphics[width=0.9\textwidth]{Abbildungen/Sequenzdiagramm.PNG}
    \caption{Sequenzdiagramm}		
    \label{fig:Sequenzdiagramm}
\end{figure}

Dem Sequenzdiagramm ist zu entnehmen, dass der Robotino nach Programmstart zyklisch Daten sendet. Die Zykluszeit beträgt 100 ms. Da die Daten über UDP an einen spezifizierten Port versendet werden ist ein Empfänger nicht zwangsweise erforderlich und die Programme können unabhängig voneinander laufen. Der Inhalt der Daten ist in Abschnitt \ref{sec:Telegramme} beschrieben. Solange der Robotino läuft werden die Daten gesendet. Dadurch kann auch ein Ausfall der Kommunikation festgestellt werden. 

Auf Seiten des Fertigungsplanungsrechners wird nur bei Bedarf ein Sendevorgang eingeleitet. Erst wenn ein Auftrag an den Robotino gesendet werden soll wird eine Schleife ausgeführt. In der Schleife wird der Auftrag zyklisch an den Robotino gesendet. Die Zykluszeit beträgt hierbei 700 ms. In Kapitel \ref{sec:StateMachineImplementierung} ist die in einer State-Machine implementierte Schleife beschrieben. Die Schleife kann durch Empfangen einer Statusänderung des Robotinos verlassen werden. Dadurch wird eine erfolgreiche Übertragung des Auftrags sichergestellt werden. Die im Auftrag befindlichen Daten sind in Abschnitt \ref{sec:Telegramme} näher erläutert.

\subsection{Telegramme}
\label{sec:Telegramme}

Die Telegramme zwischen Robotino und Fertigungsplanungsrechner bestehen aus einer festgelegten Struktur. Dabei ist sowohl die Länge der Telegramme als auch der Inhalt festgeschrieben. In beiden Telegrammen werden Double-Werte in einem Byte versendet. Jedes Byte muss daher kodiert und dekodiert werden. 

\subsubsection{Telegramm vom Fertigungsplanungsrechner zu Gewerk 2}

Der Auftrag, der an den Robotino gesendet wird, enthält 5 Bytes. Eine Aufschlüsselung ist in Tabelle \ref{tab:TelegrammZuG2} dargestellt. 

\begin{table}[!ht]
	\centering
	\begin{tabular}{|r|c|l|}
		\hline
		Byte & Inhalt	&	Beschreibung \\
		\hline
			1  & Auftragsart & 1: Transport; 2: Parken; 3: Laden  \\
			2  & Position 1  & z.B. 11 \\
			3  & Position 2  & z.B. 32 \\
		  4  & AliveStatus & 1 senden nach Alive Anfrage \\
		  5  & *Reserved   &  \\
		\hline
	\end{tabular}
	\caption{Telegramm zu Gewerk 2}
	\label{tab:TelegrammZuG2}
\end{table}

Das erste Byte klassifiziert die Auftragsart. Diese beschreibt, welche Tätigkeit der Robotino als nächstes tun soll. Auftragsart eins bedeutet, dass der Robotino einen Transportauftrag erhalten hat, also ein Werkstück von einer Position zu einer anderen Position fahren soll. Die Auftragsart 2 zeigt an, dass der Robotino auf einen Parkplatz geschickt wird. Eine Ladefahrt wird mit Auftragsart 3 gekennzeichnet. 

Im zweiten und dritten Byte sind die Positionen angegeben, an welche sich der Robotino bewegen soll. Dabei wird Position 2 nur bei Auftragsart 1 befüllt bzw. ausgewertet, da ein Parken und Laden nur eine Zielposition benötigt. Bei Auftragsart 1 jedoch zeigt Position 1 den Arbeitsplatz an, wo das Werkstück abgeholt werden soll und Position 2 den Ablageort des Werkstücks. 

Über das vierte Byte kann dem Robotino ein Status-Flag gesendet werden. Sobald der Robotino eine Anfrage zu dem Flag sendet wird eine 1 zurückgesendet. Ansonsten ist der Wert 0. Damit kann der Robotino die Funktionalität der Kommunikation validieren.

Mit dem fünften Byte wird eine einfache Erweiterung des Telegramms ermöglicht, sofern mehr Daten übertragen werden sollen. Am Anfang der Ausarbeitung des Projektes war das vierte Byte ebenfalls ein reserviertes Byte. 

\subsubsection{Telegramm von Gewerk 2 zum Fertigungsplanungsrechner}

%00: nichts; 01: unbekanntes Hindernis; 
%02: anderer Roboter; 03: direkt was gesehen?!? (flackert evtl)
%*1   X1: 1: Auf dem Weg; 2: Bin bei…; 
%X2 + X3 : Position 
%1: lebe ich noch? (Antwort senden) 0: ignorieren

\begin{table}[!ht]
	\centering
	\begin{tabular}{|r|c|l|}
		\hline
		Byte & Inhalt	&	Beschreibung \\
		\hline
			1  & Error &  Errortyp  \\
			2  & Akku  & Prozentwert als Integer \\
			3  & Hindernis & Hindernistyp \\
		  4  & *Reserved &  interner Roboterstatus \\
		  5  & AliveAbfrage  & 1: Roboterabfrage erfordert Antwort; 0: egal \\
		  6  & Status & $X_1X_2X_3$; 000: nichts zu tun \\
		  7  & Greifer   & 0:leer; 1:voll \\
		  8  & *Reserved &  \\
		  9  & *Reserved   &  \\
		\hline
	\end{tabular}
	\caption{Telegramm von Gewerk 2}
	\label{tab:TelegrammVonG2}
\end{table}

Robotererrortypen beschreiben

\begin{figure}[htb]
    \centering
    \includegraphics[width=0.9\textwidth]{Abbildungen/Positionskodierung.PNG}
    \caption{Positionskodierung}		
    \label{fig:Positionskodierung}
\end{figure}

Positionskodierung

\inlinetodo{Telegram Gewerk 2}

\inlinetodo{Port vereinbarung}

\subsection{Fehlertypen und Behebung}
\label{sec:Error}

\subsubsection{Fehlertyp 1}
\subsubsection{Fehlertyp 2}
\subsubsection{Fehlertyp 3}

% !TEX root = ../VPJ.tex

\chapter{Implementierung}
\label{sec:Implementierung}

Implementierung

\section{Programmstruktur}

\begin{figure}[htb]
    \centering
    \includegraphics[width=0.9\textwidth]{Abbildungen/Klassendiagramm.PNG}
    \caption{Klassendiagramm}		
    \label{fig:Klassendiagramm}
\end{figure}



\section{Databasehandler}
\section{UDP-Handler} 
\label{sec:UdpHandler}
\section{Robotino} 



\section{Auftrag, Prozess, Prozessschritt} 
\label{sec:AuftragProzessSchritt}
\subsection{Auftrag}
\label{sec:Auftrag}
\subsection{Prozess}
\label{sec:Prozess}
\subsection{Prozessschritt}
\label{sec:Prozessschritt}


\section{Fertigungsplanung}

Die Klasse Fertigungsplanung enthält alle verfügbaren Roboter, die Datenbankschnittstelle, die abzuarbeitenden Aufträge (vgl. Abschnitt \ref{sec:AuftragProzessSchritt}), die Fertigungsstraße und eine State Machine zur Verteilung von Roboteraufträgen. Außerdem wird auf Events reagiert, die von der Visualisierung oder anderen Quellen kommen. Zu den Events gehören Statusänderungen des Roboters, hinzufügen von Aufträgen oder Prozessen und Funktionen, die über die Hard-Code Area (vgl. Kapitel \ref{sec:HardCode}) angestoßen werden können. 
In der Fertigungsplanungsklasse werden somit alle Funktionalitäten gebündelt, die nicht direkt mit der Visualisierung zusammenhängen. 

\subsection{Initialisierung}
\label{sec:Fertigunginit}

Um die Fertigungsplanung zu initialisieren wird zunächst in der Funktion InitProzesses() ein Grundzustand hergestellt, in der mindestens 5 Prozesse vorhanden sind. Dazu werden zunächst über einen Datenbankaufruf alle dort verfügbaren Prozesse in eine Liste geschrieben. 

Sollte die Liste weniger als 5 Elemente beinhalten, also in der Datenbank zu Programmstart weniger als 5 Prozesse vorhanden sein, so werden 5 neue zuvor definierte Prozesse (Abschnitt \ref{sec:Standardprozesse}) angelegt und die Liste so überschrieben. 

Die erzeugten Prozesse werden abschließend in die Datenbank geschrieben. Außerdem erhalten alle Prozesse gemäß ihrer Schritte und Dauer in der Visualisierung einen entsprechenden Tooltip (vgl. Kapitel \ref{sec:tooltips}). 

Die Prozessliste kann über die Funktion AddProzess() dynamisch während das Programm läuft erweitert werden. Dies geschieht bei absenden eines eingestelltes Prozesses in der Visualisierung (\ref{sec:Prozesseingabe}). 

\subsubsection{Aufträge hinzufügen}

Durch Einstellen und Absenden eines Auftrags in der Visualisierung (\ref{sec:Auftragseingabe}) wird die Funktion AddAuftrag() aufgerufen. In dieser Funktion wird zunächst ein leeres Auftragsobjekt erzeugt und diesem die nächst höhere Auftragsnummer aus der Datenbank zugeordnet.  

Für jeden Prozesseintrag aus der Visualisierung wird nun die angegebene Anzahl des jeweiligen Prozesses eine Prozesskopie erzeugt. Den so kopierten Prozessen werden alle zugehörigen Prozessschritte als Kopie angehängt, um eigenständige Aufträge zu erhalten. Zuletzt werden die Prozesse und Aufträge miteinander verknüpft (siehe Listing \ref{lst:Auftragserzeugung}).

\begin{lstlisting}[frame=single, breaklines=true, numbers=left, stepnumber=2, firstnumber=1, numberstyle = \tiny, caption=Auftragserzeugung und Verknüpfungen,label=lst:Auftragserzeugung]
Auftrag *A = new Auftrag();

for (int prozesscounter = 0; prozesscounter < 5; prozesscounter++) //5 Prozess UI Elemente
{
    for (int i = 0; i < Avalues[prozesscounter]; i++) //Anzahl der eingegebenen Prozesszahl
    {
        Prozess *p = new Prozess();
        [...]  // Prozess + Prozessschritte kopieren
        A->Prozesse.append(p);
        QObject::connect(p, &Prozess::FortschrittUpdated, A, &Auftrag::UpdateFortschritt);
        QObject::connect(p, &Prozess::ProzessFinished, this, &Fertigungsplanung::StationReady);
    }
}
\end{lstlisting}

Abschließend wird der Auftrag der Auftragsliste angehängt, an die Datenbank gesendet und ein Auftragsitem in der Visualisierung erzeugt (Kapitel \ref{sec:Auftragsuebersicht}). 

\subsubsection{Standardprozesse}
\label{sec:Standardprozesse}

Folgend werden die in Abschnitt \ref{sec:Fertigunginit} erwähnten Standardprozesse näher beleuchtet. 

Um ein möglichst ausgeglichenes Fertigungsbild zu schaffen wurden die fünf Standardprozesse so untereinander abgestimmt, dass diese unterschiedlichen Kriterien genügen. 

Zunächst wurde darauf geachtet, dass die Prozesse unterschiedlich viele Prozessschritte beinhalten. Abbildung \ref{fig:Fahrweganalyse} kann entnommen werden, dass die Prozesse zwischen zwei und sechs einzelne Prozessschritte erfordern. In der Abbildung ist weiterhin dargestellt, welche Station als Start und Ziel je Prozessschritt angefahren wird. 

\begin{figure}[htb]
    \centering
    \includegraphics[width=0.9\textwidth]{Abbildungen/Fahrweganalyse.PNG}
    \caption{Diagramm Fahrweganalyse}		
    \label{fig:Fahrweganalyse}
\end{figure}

Zu erkennen ist, dass jeder Prozess bei Station 1 beginnt und in Station 8 endet. Die Einzelnen Prozesse (siehe farbliche Unterscheidung) sind dabei von unten nach oben zu lesen, beginnend mit Prozessschritt 1. 

Aus dem Diagramm können weiterhin potentielle Kollisionsherde und "`Staugefahren"' erkannt werden. Die Standardprozesse wurde so angepasst, dass alle Stationen möglichst gleichmäßig häufig durchfahren werden. Weiterhin kann aus dem Diagramm abgelesen werden, ob durch ungünstige Auftragsplanung ein Deadlock angefahren werden kann, also alle Aufträge irgendwann durch andere Aufträge blockiert sein können. Da keine Schleifen zwischen den Stationen entstanden sind kann gewährleistet werden, dass alle Aufträge in jeder Belegung nach einiger Zeit abgearbeitet werden können, wenn die Werkstücke in Station 8 entnommen werden.

Das Diagramm in Abbildung \ref{fig:Stationsanalyse} zeigt eine genauere Analyse der Stationen. Angenommen wird, dass jeder der Standardprozesse genau ein mal abgearbeitet wird. Auf der linken y-Achse, die zu den bunten Säulen gehört, ist die Belegungsdauer jeder Station in Sekunden angegeben. Die rechte, den hellgelben Säulen zugeordnete y-Achse bezeichnet die Anzahl der Anfahrten an eine Station. 

\begin{figure}[htb]
    \centering
    \includegraphics[width=0.9\textwidth]{Abbildungen/Stationsanalyse.PNG}
    \caption{Diagramm Stationsanalyse}		
    \label{fig:Stationsanalyse}
\end{figure}

Es ist zu erkennen, dass Station 1 und 8 genau fünf mal angefahren werden, also je Prozess ein mal. Die weiteren Stationen 2 bis 7 werden zwei- bis dreimal angefahren. So wird sichergestellt, dass die Stationen möglichst gleichmäßig belastet werden. 

Die farblich unterschiedlichen Säulen zeigen die kürzeste Belegungsdauer der Arbeitsplätze an. Die Balken ergeben sich durch die Dauer eines einzelnen Prozessschrittes an einem Arbeitsplatz. Es  Belegungsdauer für alle Prozesse zwischen 35 und 75 Sekunden je Station liegt. Die 2 Sekunden Bearbeitungsdauer je Prozess an Station 8 sind dabei vernachlässigt worden und gewährleisten nur ein sicheres Herausnehmen der Werkstücke durch den Menschen, ohne dass sich der Roboter noch im Entladezyklus befindet. 

Als Besonderheit ist an Station 5 zu erkennen, dass es drei Anfahrten, aber nur zwei verschiedene Prozesse gibt. Dies liegt darin begründet, dass Prozess 4 Station 5 zwei mal anfährt mit verschieden langen Zeiten. Dadurch ist die Gesamtzeit in Station 5 auch am höchsten, gehört jedoch zu großen Anteilen zum selben Prozess. Station 2 hat die kürzesten Bearbeitungszeiten je Prozess, da diese Station von drei verschiedenen Prozessen angefahren wird. Gegensätzlich dazu ist die Bearbeitungszeit an Station 23 von Prozess 3 mit 45 Sekunden sehr hoch angesetzt, da dieser Prozess keine weitere Station anfährt. 

Es wurde versucht, die durchschnittliche Belegungsdauer der Stationen anzugleichen, um alle Arbeitsplätze möglichst homogen auszulasten. Zu erkennen ist, dass die minimale

Als Nebeneffekt ergibt sich, dass die Roboter möglichst ausgeglichen über die gesamte Fertigungsstraße fahren und die Gefahr von Staus oder Kollisionen dadurch verringert wird. 

\subsection{Roboter Statusänderungen}

Wenn eine der in der Fertigungsplanung enthalten Roboterklassen ein Event feuert, welches eine Statusänderung beinhaltet (vgl. Kapitel \ref{sec:UdpHandler}), wird über Funktionen in der Fertigungsplanung auf diese reagiert. So wird der Greifer oder das Statusflag eines Roboters bei einer Änderung aktualisiert. Beim Empfanges eines Errorevents des Roboter wird je nach Errortyp eine Meldung ausgegeben und individuell auf die Robotererror reagiert. 

Weiterhin wird durch Ablauf des Timers eines Arbeitsplatzes dieser in der Fertigungsstraße aktualisiert und auf bereit geschaltet. 


\subsubsection{Roboterstatus geändert}

Bei Änderung des Roboterstatus wird die Funktion OnRobotStatusChanged() aufgerufen. Dabei wird zunächst der aufrufende Roboter ausgewählt und zwischengespeichert. Sollte sich der Status des Roboter auf eins geändert haben, also der Roboter auf dem Weg zu etwas ist, so wird unterschieden ob sich der Roboter auf einem Parkplatz oder einem Arbeitsplatz befunden hat.

Von einem Parkplatz (auch Ladestation) kommend wird dieser in der Datenbank und der Visualisierung wieder freigegeben. Wenn der Roboter zuletzt an einer Station war, so wird diese jetzt für andere Roboter freigegeben und der Roboter wird in der Datenbank aus dem Arbeitsplatz gelöscht. 

\subsubsection{Greiferstatus geändert}

Die Prozessaktualisierungen basieren auf der Änderung des Greifers des Roboters. Es wird vorausgesetzt, dass sich der Greifer zu keinem Zeitpunkt willkürlich öffnet oder schließt, sondern nur bei Aufnahme oder Abgabe eines Werkstücks an einer Station. 

Das Schließen des Greifers eines Roboters bedeutet somit immer, dass ein Werkstück aufgenommen wurde. Somit wird in der Datenbank der Arbeitsplatz, an dem sich der Roboter befindet freigegeben und kann für andere Aufträge genutzt werden. Der Arbeitsplatz wird in der Visualisierung und der Datenbank freigegeben. Zuletzt wird der Tooltip des Arbeitsplatzes (vgl. Kapitel \ref{sec:tooltips}) auf einen leeren String gesetzt. 

Beim Öffnen des Greifers wird ein Werkstück in einen Arbeitsplatz einer Station gelegt. Daraufhin wird die Zuordnung des Werkstücks zu dem entsprechenden Roboter gelöscht. Mit der Werkstücks-ID kann nun aus der Auftragsliste der zugehörige Prozess identifiziert werden. Der Timer des so gefundenen Prozesses wird jetzt gestartet, somit läuft die Bearbeitung des Werkstücks an dem Arbeitsplatz. In der Visualisierung und Datenbank wird der Arbeitsplatz von reserviert auf belegt aktualisiert.

Weiterhin wird das Werkstück, dass dem Roboter zu diesem Zeitpunkt zugeordnet ist an der Station an der sich der Roboter zuletzt aufgehalten hat gelöscht und der zugehörige Arbeitsplatz wird freigegeben. 

\subsubsection{Roboter Error geändert}

Bei Auftreten eines Errors im Roboter wird je nach Typ eine Error-Meldung ins Log geschrieben. 

Die Error-Meldungen des Roboters wurden bewusst nicht automatisiert behandelt, da dem Benutzer die Kontrolle über das Gesamtsystem behalten sollte. Somit muss je nach Fehlertyp eine unterschiedlicher Lösungsweg verfolgt werden. 

Die auftretenden Fehler und der Lösungsweg ist in Kapitel \ref{sec:Error} zu finden. 

\subsection{Hard-Code Funktionen}

Die mit der Hard-Code Area erzeugten Befehle aus der Visualisierung (vgl. Kapitel \ref{sec:HardCode}) werden in verschiedenen Funktionen verarbeitet. So wird beispielsweise der Parkplatz oder eine Station in Datenbank und Visualisierung freigegeben, ein Roboter defekt geschaltet oder ein Arbeitsplatz als defekt markiert oder repariert. 

\subsection{Zustandsdiagramm}

Das Zustandsdiagramm in Abbildung \ref{fig:Auftragsplanung} beschreibt die Logik der Auftragsvergabe an die Roboter. Das Diagramm wurde anschließend als State Machine (vgl. \ref{sec:StateMachines}) implementiert. 

\begin{figure}[htb]
    \centering
    \includegraphics[width=0.8\textwidth]{Abbildungen/Auftragsplanung_rotated.PNG}
    \caption{Zustandsdiagramm der Auftragsplanung}		
    \label{fig:Auftragsplanung}
\end{figure}

In der Abbildung sind die Zustände mit Namen dargestellt. Alle Zustände haben eine komplementäre Funktion, die nach einem Zustandswechsel als Entry-Action aufgerufen wird. Eine genaue Funktionsbeschreibung der Zustände ist in Kapitel \ref{sec:StateMachineImplementierung} zu finden. Um zwischen den Zuständen zu wechseln werden die mit Pfeilen eingezeichneten Transitionen verwendet. Die Pfeilbeschriftung beinhaltet dabei das Signal, welches für den Zustandswechsel genutzt wird. Um das Diagramm übersichtlich zu gestalten, wurden die Timeout-Transitionen nicht vollständig eingezeichnet, sondern enden in einem kreisförmigen Xa-Zustand. Dieser führt direkt zu dem XM-Zustand oben links und leitet in den Zustand Timeout weiter. 

\subsection{State Machine Implementierung}
\label{sec:StateMachineImplementierung}

Im Quellcode werden zunächst alle Zustände erzeugt und der State Machine hinzugefügt. Anschließend werden alle Transitionen zwischen den Zuständen erzeugt und der State Machine hinzugefügt. Hierbei werden definierte Signale genutzt, die zuvor in der Header Datei deklariert wurden. In der Main-Funktion werden abschließend die Zustände mit den entsprechenden zugehörigen Funktionen verknüpft (vgl. \ref{sec:Main}). 

Um die State Machine zu starten wird, nachdem der Initialzustand bekannt gemacht wurde, die Start-Funktion aufgerufen. Eine verkürzte Darstellung der Aufrufe ist in Listing \ref{lst:StartStateMachine} abgebildet. 

\begin{lstlisting}[frame=single, breaklines=true, numbers=left, stepnumber=2, firstnumber=1, numberstyle = \tiny, caption=Start und Initialisierung der State Machine,label=lst:StartStateMachine]
    QStateMachine StateMachine;
    QState *StateCheckRobots = new QState;

    StateMachine.addState(StateCheckRobots);
    StateCheckRobots->addTransition(this, SIGNAL(nextState()), StateCheckAkku);

    StateMachine.setInitialState(StateCheckRobots);
    StateMachine.start();

    QObject::connect(StateCheckRobots, &QState::entered, this, &Fertigungsplanung::CheckRobots);
\end{lstlisting}

\subsubsection{Zustand CheckRobots}

Im Initialzustand der State Machine, CheckRobots, werden alle vier Roboter nacheinander auf Bereitschaft überprüft. Die Reihenfolge der zu überprüfenden Roboter wird dabei durch einen hierfür entwickelten Shuffle-Algorithmus (\ref{sec:shuffle}) bestimmt. 

Ein Roboter gilt als bereit, wenn er als Status 0 zurückgibt, nicht als defekt markiert ist und als lebend gilt. Sobald ein geprüfter Roboter alle drei Kriterien erfüllt wird dieser ausgewählt und in den weiteren Zuständen, gefolgt von CheckAkku, genutzt. Die anderen Roboter werden bis zum nächsten Aufruf des Zustands CheckRobots vernachlässigt.

Sollte keiner der vier Roboter alle drei Kriterien erfüllen wird in den Zustand CheckParkplatzRobots gewechselt.

\subsubsection{Zustand CheckParkplatzRobots}

Wie im Zustand CheckRobots werden alle vier Roboter anhand von Kriterien auf Bereitschaft überprüft. Dazu wird ebenfalls die Reihenfolge zufällig bestimmt.

Anstatt den Roboterstatus wie beim Zustand CheckRobots auf 0 zu prüfen, wird ein Status zwischen 201 und 204 erwartet. Diese vier Stati repräsentieren die vier Parkplätze auf denen sich ein Roboter befinden kann. 

Durch die Aufteilung der Zustände CheckRobots und CheckParkplatzRobots erfolgt eine Priorisierung der Roboter, die sich bereits in den Stationen befinden und einen Auftrag fertig abgearbeitet haben. Nur wenn kein Roboter in den Stationen verfügbar ist wird auf die auf den Parkplatz befindlichen Roboter zurückgegriffen. Dadurch werden Situationen vermindert, in denen sich die Roboter bei der Stationsarbeit gegenseitig behindern.

Sollte keiner der auf den Parkplatz befindlichen Roboter alle weiteren Kriterien erfüllen wird in den Zustand Wait gewechselt.

\subsubsection{Zustand Wait}

Da in Qt aufgrund der eventbasierten Programmierung (vgl. \ref{sec:Eventbasiert}) keine Endlosschleifen vorgesehen sind wurde im Zustand "`Wait"' mittels Timer die Zustandsmaschine künstlich verlangsamt. Ein Entfernen des Zustands führt zum Programmabsturz, da in kürzester Zeit die interne Eventliste voll geschrieben wird. Auf externe Events wie Mauseingaben oder Ereignisse vom Betriebssystem könnte somit nicht mehr reagiert werden. 

Durch den Zustand Wait können andere Programmevents und Ereignisse des Betriebssystems abgearbeitet werden, während 2000 ms auf das timeout()-Signal des Timers gewartet wird. 

Nach Ablauf des Timers wird in den Initialzustand CheckRobots gewechselt.

\subsubsection{Zustand CheckAkku}

Der Zustand CheckAkku gewährleistet einen Tiefentenladungsschutz der Roboter. Sobald ein Roboter in den ersten beiden Zuständen ausgewählt wurde wird überprüft, ob sein Akkustand größer gleich 25 Prozent ist. Es wird davon ausgegangen, dass kein Auftrag plus eine Fahrt zum Laden mehr als 25 Prozent Akku verbraucht. Der Wert wurde empirisch ermittelt. 

Sollte der Akkustand unter 25 Prozent liegen, wird der Roboter über Zustand SendLoadingTask zu einer freien Ladestation geschickt. Wenn keine Ladestation frei ist wird in den Zustand Wait gewechselt. 
Bei einem Akkustand von über 25 Prozent wird der Zustand CheckAuftrag aufgerufen. 

\subsubsection{Zustände SendLoadingTask, SendParkingTask, SendPositionTask}

Um bei Kommunikation mit den Robotern in den Zuständen SendLoadingTask, SendPositionTask und SendParkingTask sicherzustellen, dass die Nachricht empfangen wurde, wird diese mehrfach gesendet, bis die Übertragung erfolgreich war (siehe Kapitel \ref{sec:sequenzdiagram}). Es wird nach jedem Sendevorgang ein kurzer Wartezustand aufgerufen, der, gleich dem Zustand Wait, einen Programmabsturz verhindert. Mit dem individuellen Wartezustand kann ebenfalls das bei aktuell 700 ms angesetzte Sendeintervall eingestellt werden. 

Aus allen den drei Send......Task Zuständen und ihren zugehörigen Wartezuständen kann über ein Ablauf eines Timers in den Zustand Timeout gewechselt werden. Der Timer wird beim erstmaligen Eintritt in einen der Zustände Send......Task gestartet und läuft nach 25 Sekunden ab. Dieser Timer gewährt eine Absturzsicherheit des Programms bei Roboterausfall oder anderweitigem Fehlschlagen der Kommunikation. 

In allen drei Send.....Task Zuständen wird ein erfolgreicher Sendevorgang damit erkannt, dass der Roboterstatus nicht mehr 0 ist, und der Roboter sich nicht mehr auf einem Parkplatz befindet. Solange die beiden Bedingungen nicht erfüllt sind wird weiter an den Roboter gesendet. 

Wenn im Zustand SendLoadingTask die Nachricht erfolgreich versendet wurde, wird der Timer für das Timeout gestoppt, und die benötigte Ladestation in der Visualisierung und in der Datenbank reserviert. Abschließend wird in den Initialzustand gewechselt.

Der Zustand SendParkingTask verhält sich Analog dem SendLoadingTask, nur wird bei erfolgreicher Nachrichtenversendung der benötigte Parkplatz reserviert.

Bei dem Zustand SendPositionTask, in dem ein Auftrag aus der Planung an den Roboter verschickt wurde, werden nach erfolgreichem Nachrichtenversand mehrere Aktionen durchgeführt. Zunächst wird der Timeout-Timer gestoppt. Sollte der Auftrag den ersten Prozessschritt, also das Abholen eines Werkstücks im Lager an Station 1, beinhalten wird der Schreibkopf des RFID neu beschrieben. In der Datenbank wird außerdem dem Start- und Zielarbeitsplatz der genutzte Roboter zugeordnet. Am Startarbeitsplatz wird das vorhandene Werkstück in der Datenbank entfernt und dem Zielarbeitsplatz zugeordnet. Die beiden Stationen werden außerdem Reserviert um zu verhindern, dass ein anderer Roboter an die zu bearbeitenden Stationen geschickt wird. Somit ist eine erste Kollisionsvermeidung implementiert. Zuletzt wird in der Datenbank das genutzte Werkstück dem Roboter zugewiesen. Auch in der Visualisierung werden Start- und Zielarbeitsplatz reserviert. Schlussendlich wird in der Visualisierung der Tooltip (vgl. Kapitel \ref{sec:tooltips}) des Start- und Zielarbeitsplatz aktualisiert.

\subsubsection{Zustand CheckAuftrag}

Der Zustand CheckAuftrag enthält die Vergabe der Aufträge an die Roboter. 
Dazu wird zunächst in der Auftragsliste jeder Auftrag auf unfertige Prozesse untersucht und diese folgend an eine neu erzeugte Prozessliste angehängt (vgl. Abschnitt \ref{sec:AuftragProzessSchritt}). Jeder Prozess der so erzeugten Prozessliste wird anschließend auf die Möglichkeit der Bearbeitung geprüft. Sofern eine der folgenden Prüfungen fehlschlägt wird der nächste Prozess geprüft.

Dazu wird zunächst überprüft ob die Auftragsplanung pausiert ist oder sich der Prozess schon in Bearbeitung befindet. 

Es wird geprüft ob die Start und Zielstation nicht durch einen Roboter reserviert ist.

Sollte der nächste Prozessschritt der bearbeitet werden muss an der Lagerstation (Station 1) sein, so wird gewährleistet, dass sich ein Werkstück im Lager befindet und dieses ausgewählt. 

Es wird geprüft ob eine der beiden Zielarbeitsplätze der Zielstation frei ist, also weder reserviert noch defekt und anschließend ausgewählt. 

Wenn alle Vorbedingungen für einen der Prozesse eintreffen wird der Zustand direkt verlassen und der Auftrag wird an den ausgewählten Roboter im Zustand SendPositionTask gesendet. 

Sollten alle Prozesse zu keinem Sendevorgang geführt haben, das heißt es ist entweder kein Auftrag vorhanden oder oder das senden wurde blockiert, so wird der aktuell ausgewählte Roboter zum Parken geschickt (Zustand SendParkingTask), sofern er sich noch nicht auf einem Parkplatz befindet. Sonst wird in den Initialzustand gewechselt.

\subsection{Shuffle-Algorithmus}
\label{sec:shuffle}

Der Shuffle-Algorithmus dient dazu, die Roboter in zufälliger Reihenfolge abzuarbeiten, um eventuelle Dead-Locks zu vermeiden. Ein solcher Dead-Lock könnte zum Beispiel eintreten, wenn Roboter 1 und 2 an der Ladestation sind, Roboter 3 zum laden geschickt werden müsste, jedoch warten muss. Roboter 4 würde dabei, selbst wenn er bereit wäre einen Auftrag abzuarbeiten nie überprüft werden.

Zurückgegeben wird eine Liste in der die Zahlen 1 bis 4 in zufälliger Reihenfolge vorliegen. Der Algorithmus ist optimiert gegenüber der Anzahl an Systemaufrufen für eine Zufallszahl. Andernfalls wäre es einfacher solange eine Zufallszahl zurückgeben zu lassen,  wie die Liste diese noch nicht enthält.

Um die Roboterreihenfolge festzulegen werden vier Schritte durchgeführt. Im ersten Schritt wird eine leere Liste erzeugt und eine zufällige Zahl zwischen 1 und 4 an die erste Stelle geschrieben (Listing \ref{lst:shuffle} Zeile 1-3). 

\begin{lstlisting}[frame=single, breaklines=true, numbers=left, stepnumber=2, firstnumber=1, numberstyle = \tiny, caption=Shuffle-Algorithmus,label=lst:shuffle]
    QList<int> liste;
    int i = QRandomGenerator::global()->bounded(1, 5);
    liste.append(i);
    i = QRandomGenerator::global()->bounded(1, 4);
    (liste.contains(i)) ? liste.append(i+1) : liste.append(i);
    i = QRandomGenerator::global()->bounded(1, 3);
    if (liste.contains(i))
    {
        if (liste.contains(i+1))
        {
            liste.append(i+2);
        }
        else
        {
            liste.append(i+1);
        }
    }
    else
    {
        liste.append(i);
    }

    if (!liste.contains(1)) {liste.append(1);}
    else if (!liste.contains(2)) {liste.append(2);}
    else if (!liste.contains(3)) {liste.append(3);}
    else if (!liste.contains(4)) {liste.append(4);}
    return liste;
\end{lstlisting}

Im zweiten Schritt wird eine Zufallszahl zwischen 1 und 3, wenn sie noch nicht enthalten ist, in die Liste geschrieben oder um eins inkrementiert und in die Liste geschrieben (Listing \ref{lst:shuffle} Zeile 4f). 

Um die nächste Zahl hinzuzufügen wird im dritten Schritt (Listing \ref{lst:shuffle} Zeile 6-21) eine Zufallszahl zwischen 1 und 2 erzeugt und in die Liste geschrieben, sollte sie nicht vorhanden sein. Wenn sie schon existiert wird sie um eins inkrementiert und erneut überprüft ob die inkrementierte Zahl in der Liste ist. Aufgrund der Maximalgröße von 4 kann der Algorithmus so alle Fälle abdecken.
 
Zuletzt wird die noch fehlende Zahl der Liste ergänzt (Listing \ref{lst:shuffle} Zeile 22ff) und die Liste zurückgegeben. 

\section{Mainwindow}

\section{Main}
\label{sec:Main}
\inlinetodo{StateMachine Connections beschreiben / erwähnen - genau eine funktion zu einem State}
\section{Weitere Klassen}



% !TEX root = ../VPJ.tex

\chapter{Visualisierung}
\label{sec:Visualisierung}

Die Auftragsplanungssoftware besitzt neben der Logik eine Visualisierung. Diese wurde mit den von Qt bereitgestellten Widgets grafisch programmiert und in der Funktion mainwindow verknüpft und organisiert. 

\section{Design}

Die Leitlinie des Grafikdesigns der HAW war:

\begin{quote}
"`Fokussiert. Direkt. Verständlich.
Und vor allem sympathisch."'	

"`Idee der neuen visuellen Sprache ist die Reduktion auf das
Wesentliche. Ziel soll es sein, einfach und klar die wichtigste
Aussage zu kommunizieren."'
\cite[4]{corporate}
\end{quote}

Beim Design der Software wurde versucht diese Aspekte wieder aufzugreifen. Die Visualisierung sollte ebenso direkt und verständlich sein. Vor allem die Reduktion auf das Wesentliche lag besonders im Fokus. 

Um die Visualisierung im Stil der HAW zu gestalten wurde das Corporate Design bei Farb und Schriftwahl verwendet. Somit sind alle Schriftarten in der Visualisierung Open Sans. Als Farben wurden die vorgegebenen Farbstile "`HAW Hauptblau"' als Umrahmung und Abtrennung aller Bereiche gewählt. Das "`HAW Hellblau"' ist innerhalb der verschiedenen Bereiche als Trennung der Bedienelemente zu sehen. Da dieses blasser ist können alle Bedienelemente gut erkannt und verwendet werden. 

\begin{figure}[htb]
    \centering
    \includegraphics[width=0.3\textwidth]{Abbildungen/VPJ.png}
    \caption{VPJ-Logo}		
    \label{fig:VPJLogo}
\end{figure}

Das entworfene VPJ-Logo (Abb. \ref{fig:VPJLogo}) welches in Taskmanager, Taskleiste, an oberer Bildschirmseite und zentral in der Mitte der Visualisierung zu finden ist, unterliegt ebenfalls den Richtlinien, die im Corporate Design (vgl. \cite{corporate}) festgehalten sind. 

\section{Struktur}

Die Visualisierung lässt sich in insgesamt 6 Bereiche unterteilen. Die Bereiche sind in Abbildung \ref{fig:GesamtprogrammBereiche} dargestellt. Bereich I oben links zeigt die Fertigungsstraße mitsamt Robotern, Parkplätzen und Stationen. Dieser Bereich wird folgend Live-View genannt. 

\begin{figure}[htb]
    \centering
    \includegraphics[width=0.8\textwidth]{Abbildungen/GesamtprogrammBereiche.png}
    \caption{Übersicht über die Bereiche innerhalb der Visualisierung}		
    \label{fig:GesamtprogrammBereiche}
\end{figure}


In Bereich II, der Auftragsübersicht, können die Fortschritte der laufenden Aufträge und Prozesse überwacht werden.

Innerhalb des Bereich III können über Tabs vier weitere Fenster aufgerufen werden. Dabei dienen zwei der Anzeige spezifizierter Informationen, wohingegen die anderen beiden Tabs der Kontrolle dienen. Die Kombination aller vier Tabs wird folgend Tab-View genannt.

Bereich IV zeigt den Status der Robotinos an. 

Die beiden Bereiche V und VI werden ausschließlich für die Erzeugung neuer Aufträge oder Prozesse verwendet. 

In Abbildung \ref{fig:Gesamtprogramm} ist eine Übersicht der Visualisierung im laufenden Betrieb dargestellt. 

\begin{figure}[htb]
    \centering
    \includegraphics[width=1\textwidth]{Abbildungen/Gesamtprogramm.png}
    \caption{Übersicht Visualisierung}		
    \label{fig:Gesamtprogramm}
\end{figure}

In den folgenden Kapiteln werden die einzelnen Bereiche der Abbildung näher beleuchtet und ihre Implementierung in der MainWindow Funktion erläutert. 

\section{Live-View}

\begin{figure}[htb]
    \centering
    \includegraphics[width=0.9\textwidth]{Abbildungen/LiveView.png}
    \caption{Visualisierung Live-View}		
    \label{fig:LiveView}
\end{figure}

\subsection{Parkplätze und Ladestationen}

\begin{figure}[htb]
    \centering
    \includegraphics[width=0.4\textwidth]{Abbildungen/Parkplatz.png}
    \caption{Visualisierung Parkplatz}		
    \label{fig:Parkplatz}
\end{figure}

\begin{figure}[htb]
    \centering
    \includegraphics[width=0.25\textwidth]{Abbildungen/Laden.png}
    \caption{Visualisierung Ladestation}		
    \label{fig:Ladestation}
\end{figure}

\subsection{Stationen}

\begin{figure}[htb]
    \centering
    \includegraphics[width=0.18\textwidth]{Abbildungen/Station.png}
    \caption{Station mit verschiedenen Arbeitsplatzstati}		
    \label{fig:Station}
\end{figure}

\subsection{Robotino}

\begin{figure}[htb]
    \centering
    \includegraphics[width=0.13\textwidth]{Abbildungen/RobotinoGoffen.png}
    \includegraphics[width=0.13\textwidth]{Abbildungen/RobotinoGZu.png}
    \includegraphics[width=0.13\textwidth]{Abbildungen/RobotinoDefect.png}
    \caption{Roboter in verschiedenen Stati (vlnr: Greifer offen, Greifer geschlossen, Defekt)}		
    \label{fig:Robotino}
\end{figure}

\section{Auftragsübersicht}
\label{sec:Auftragsuebersicht}

\begin{figure}[htb]
    \centering
    \includegraphics[width=0.4\textwidth]{Abbildungen/Auftragsfortschritt.png}
    \caption{Auftragsfortschritt}		
    \label{fig:Auftragsfortschritt}
\end{figure}

\begin{figure}[htb]
    \centering
    \includegraphics[width=0.4\textwidth]{Abbildungen/Pause.png}
    \caption{Pause Button gedrückt}		
    \label{fig:Pause}
\end{figure}

\section{Tab-View}

\begin{figure}[htb]
    \centering
    \includegraphics[width=0.6\textwidth]{Abbildungen/Log.png}
    \caption{Tab: Log View}		
    \label{fig:Log}
\end{figure}

\begin{figure}[htb]
    \centering
    \includegraphics[width=0.6\textwidth]{Abbildungen/ManualControl.png}
    \caption{Tab: Manual Control}		
    \label{fig:ManualControl}
\end{figure}

\begin{figure}[htb]
    \centering
    \includegraphics[width=0.4\textwidth]{Abbildungen/TimestampsWerkstueck.png}
    \includegraphics[width=0.4\textwidth]{Abbildungen/TimestampsStation.png}
    \caption{Tab: Info Area}		
    \label{fig:InfoArea}
\end{figure}

\subsection{Hard-Code Bereich}
\label{sec:HardCode}

\begin{figure}[htb]
    \centering
    \includegraphics[width=0.6\textwidth]{Abbildungen/HardCode.png}
    \caption{Tab: Hard-Code}		
    \label{fig:HardCode}
\end{figure}

\begin{figure}[htb]
    \centering
    \includegraphics[width=1\textwidth]{Abbildungen/GesamtprogrammROT.png}
    \caption{Übersicht Visualisierung im Hard-Code Modus}		
    \label{fig:GesamtprogrammROT}
\end{figure}

\section{Roboterstatus}

\begin{figure}[htb]
    \centering
    \includegraphics[width=0.5\textwidth]{Abbildungen/Batterie.png}
    \caption{Batterie und Statusanzeige}		
    \label{fig:Batterie}
\end{figure}

\begin{figure}[htb]
    \centering
    \includegraphics[width=0.4\textwidth]{Abbildungen/BatterieAlive1.png}
    \includegraphics[width=0.4\textwidth]{Abbildungen/BatterieAlive2.png}
    \caption{Roboterstatus-LED blinkend}		
    \label{fig:Led}
\end{figure}

\section{Prozesseingabe}
\label{sec:Prozesseingabe}

\begin{figure}[htb]
    \centering
    \includegraphics[width=0.5\textwidth]{Abbildungen/Prozesseingabe.png}
    \caption{Visualisierung - Prozesseingabe}		
    \label{fig:Prozesseingabe}
\end{figure}

\section{Auftragseingabe}
\label{sec:Auftragseingabe}

\begin{figure}[htb]
    \centering
    \includegraphics[width=0.25\textwidth]{Abbildungen/Auftragsvergabe.png}
    \caption{Visualisierung - Auftragsvergabe}		
    \label{fig:Auftragsvergabe}
\end{figure}


\section{Benutzerinteraktion}

\subsection{Tooltips}
\label{sec:tooltips}
\inlinetodo {Tooltips mit Bildern her und erlaeeeren}

\section{MainWindow}

Neben allen bisher bereits beschriebenen Funktionen wird folgend eine Übersicht über die MainWindow-Funktion gegeben. 


% !TEX root = ../VPJ.tex

\chapter{Simulation}
\label{sec:Simulation}


\begin{figure}[htb]
    \centering
    \includegraphics[width=0.9\textwidth]{Abbildungen/SimulatedRobot.PNG}
    \caption{Ablaufdiagramm des simulierten Roboters}		
    \label{fig:simRobot}
\end{figure}



\chapter{Fertigungsrechner}\label{kap:Fertigungsrechner}
 Der Fertigungsrechner ist Teil des Gewerk 1 und besteht hauptsächlich aus der MySQL-Datenbank und der Fertigungsüberwachung. Die MySQL-Daten\-bank be\-in\-hal\-tet alle relevanten Daten des Fertigungsprozesses. Der Entwurf und die Implementierung der Datenbank wird in Kapitel \ref{kap:Datenbank} beschrieben. Über das lokale LAN-Netzwerk kann der Rechner der Fertigungsplanung auf den Fertigungsrechner beziehungsweise die hier abgelegte MySQL-Datenbank zugreifen. 
 Die Fertigungsüberwachung steuert die Ansteuerung und Auswertung der RFID-Schreib-Lese-Köpfe und legt die Daten der Werkstückverfolgung in der Datenbank ab. Mit den Schreib-Lese-Köpfen ist eine Verfolgung der Werkstücke im Fertigungsprozess möglich, da jeder An- und Abmeldevorgang an den Fertigungsstationen 2-7 und dem Rohteil- so wie dem Fertigteillager überwacht und gespeichert wird.  Über ein Interface und das Bussystem Profibus-DP sind die Schreib-Lese-Köpfe an den Fertigungsrechner angeschlossen. Die Fertigungsüberwachung wird in Kapitel \ref{kap:Fertigungsueberwachung} beschrieben.  In der Abbildung \ref{fig:Fertigungsrechner} ist eine Übersicht über die Komponenten des Fertigungsrechners und die Schnittstellen dargestellt.
 \begin{figure}[hbt]
	    \centering
	    \includegraphics[width=1\linewidth]{Bilder/Fertigungsrechner.png}
        \caption{Fertigungsrechner und Schnittstellen}
        \label{fig:Fertigungsrechner}
\end{figure}
 
% !TEX root = ../VPJ.tex

\chapter{Datenbank}
\label{sec:Datenbank}



\chapter{Fertigungsüberwachung}\label{Fertigungsueberwachung}
Die Fertigungsüberwachung kontrolliert die Fertigung der Werkstücke und ermöglicht eine Verfolgung der Werkstücke während des Fertigungsprozesses. Dazu nutzt die Fertigungsüberwachung RFID-Lese-/Schreibköpfe die über Profibus mit dem Fertigungsrechner verbunden sind und speichert die ermittelten Timestamps in der im Kapitel \ref{kap:Datenbank} beschriebenen Datenbank ab.
Die Fertigungsüberwachung befindet sich ebenso wie die Datenbank auf dem Fertigungsrechner und ist als Soft-SPS ausgeführt. Für die Programmierung wird CODESYS genutzt. Für den Zugriff auf die Datenbank wird wie in Abschnitt \ref{kap:DatenbankZugriff} ausgeführt das Tool SQL4Automation genutzt.

\section{Konzept}
Das Programm ist so organisiert, das verschiedene Programmteile quasi Parallel ablaufen. Dazu sind die einzelnen Programmteile in Funktionsbausteine ausgelagert. Die Instazen dieser Funktionsbausteine werden in der Main oder in den jeweiligen Instanzen aufgerufen. So lässt sich das Programm Strukturieren. In Abbildung \ref{fig:FB_Uebersicht} ist die Struktur der Funktionsbausteine und deren Instanzen dargestellt. Dabei ist immer der Instanzenname und in Klammern der Name des Funktionsbausteins angegeben. Türkis hinterlegt sind die Stellen wo ein Funktionsbaustein global instanziert wurde und innerhalb der lokalen Instanz genutzt wird.
\begin{figure}[h]
	    \centering
	    \includegraphics[width=0.7\linewidth]{Bilder/FB_Uebersicht_final.png}
        \caption{Übersicht über die Funktionsbausteine und ihre Instanzen}
        \label{fig:FB_Uebersicht}
\end{figure}
Nähere Informationen zu den einzelnen Funktionsbausteinen sind in Abschnitt \ref{kap:FBs} zu finden. Innerhalb der Funktionsbausteine ist das Programm überwiegend durch Zustandsautomaten organisiert. Dies ermöglicht ein Strukturiertes Programm und den grafischen Entwurf mithilfe eines UML-Zustandsdiagramms.

Aufgeteilt ist das Programm in drei Funktionsteile. Der eine Teil sorgt für das Auslesen und Beschreiben der RFID-Tags. Dabei wird einerseits ein Abbild der aktuell Eingelesenen Werte der RFID-Tags unter den Lesegeräten der Stationen 2 bis 8 erstellt und andererseits der Station 1 der Wert zugewiesen mit der der nächste RFID-Tag beschrieben werden soll. Dieser Teil befindet sich in den Funktionsbausteinen \glqq Aufruf\_RFID\_Channel\_2\_bis\_8\_Lesen\grqq  und \glqq Aufruf\_RFID\_Channel\_1\_Schreiben\grqq . Der zweite Teile sorgt für die Kommunikation zwischen der Soft-SPS und der Datenbank und befindet sich im global instanzierten Funktionsbaustein \glqq Datenbank\_Read\_Write\grqq . Der letzte Teil sorgt für die Kommunikation zwischen den beiden anderen Teilen. Die eingelesenen Werte der RFID-Lesegeräte und die aus Datenbankabfragen gewonnen Werte werden ausgewertet und daraus Anweisungen an die anderen Programmteile generiert. 

\section{Funktionsbausteine}\label{kap:FBs}

\subsection{FB Aufruf\_RFID\_Channel\_2\_bis\_8\_Lesen} \label{kap:FB_Lesen}
Dieser Funktionsbaustein ließt die RFID-Leseköpfe an den Stationen 2 bis 8 aus. Es handelt sich hierbei um eine Erweiterung des im Staterkit zurverfügunggestellten Funktionsbausteins \glqq Aufruf\_RFID\_Channel\_1\_bis\_8\grqq . Die Leseköpfe 2 bis 8 werden nur lesend betrieben. Der eigentliche Zugriff auf die RFID-Leseköpfe findet in den sieben Instanzen des auch mit dem Staterkit zurverfügunggestellten Funktionsbausteins \glqq FB\_RFID\_SS15\grqq  statt. Auch die Initialisierung findet in diesen Instanzen statt und muss über die entsprechende Eingangsvariable gestartet werden .

Die Instanzen des Funktionsbausteins \glqq FB\_RFID\_SS15\grqq  werden in jedem Zyklus aufgerufen. Zu beginn des Funktionsbausteins befindet sich ein Zustandsautomat, der die Initialisierung der Instanzen des  \glqq FB\_RFID\_SS15\grqq  durchführt und den Lesebetrieb festlegt. Um den Code kürzer und Übersichtlicher zu gestalten befindet sich der Zustandsautomat in einer for-Schleife, wobei jeder Schleifendurchlauf einem Lesekopf entspricht. Die entsprechenden Variablen für die Initialisierung befinden sich dazu in einem Array. Der Zustandsautomat für die Initialisierung ist in Abbildung \ref{fig:init_Leseköpfe} dargestellt.

\begin{figure}[h]
	    \centering
	    \includegraphics[width=0.7\linewidth]{Bilder/Zustandsdiagramme/RFID_lesen.png}
        \caption{UML-Zustandsdiagramm der Initilisierung der RFID-Leseköpfe 2 bis 8}
        \label{fig:init_Leseköpfe}
\end{figure}

Ist die Initalisierung abgeschlossen werden die RFID-Leseköpfe kontinuierlich ausgelesen. Ist ein RFID-Tag unter dem Lesegerät wird der Wert des Tags in einem Array gespeichert. Dabei entspricht die Stelle im Array der Station, wo der Tag eingelesen wurde. Ist kein Werkstück unter dem RFID-Lesegerät wird in das Array eine 0 geschrieben. Da der Wert den die RFID-Lesegeräte zurückliefern immer dem des zuletzt eingelesenen RFID-Tags entspricht wird mithilfe des Parameters TFR überprüft ob aktuell ein RFID-Tag unter dem Lesegerät ist. Der TFR-Parameter gibt an ob sich ein RFID-Tag innerhalb des Lesebereichs des RFID-Lesekopf befindet. Befindet sich kein RFDI-Tag unter dem Lesekopf wird der Wert auf 0 gesetzt. Befindet sich ein RFID-Tag unter dem Lesegerät wird gewartet bis der Lesekopf den RFID-Tag ausgelesen hat und der Wert ins Array geschrieben. 

Da der RFID-Tag als einzelne Bytes ausgelesen wird müssen diese noch in einen Integer-Wert umgerechnet werden. Dies übernimt die selbstgeschriebene Funktion ByteToInt(), die als Eingangsvariable ein Byte-Array und als Ausgangsvariable einen Integervariable hat.
\subsection{FB Aufruf\_RFID\_Channel\_1\_Schreiben}
Dieser Funktionsbaustein organisiert das beschreiben der RFID-Tags an Station 1. Im Aufbau ist er dem in Abschnitt \ref{kap:FB_Lesen} beschriebenen Funktionsbasutein ähnlich. Allerdings befindet sich hier nur eine Instanz des Funktionsbausteins \glqq FB\_RFID\_SS15\grqq . Der Zustandsautomat entspricht dem in Abbildung \ref{fig:init_Leseköpfe}. Lediglich in Zustand 2 wird der RFID-Lese-/Schreibkopf nicht als Lesekopf sondern als Schreibkopf eingerichtet. 

Über die Rückgabewerte Busy und Done wird detektiert wenn ein RFID-Tag neu beschrieben wurde. Wird ein Schreibvorgang erkannt wird die Globale Variable \glqq WRITE\_RFDI\_DONE\grqq gesetzt. Diese Variable wird im Funktionsbaustein \glqq Timestamp\_anlegen\grqq aus Abschnitt \ref{kap:FB_Timestamp} beim Suchen nach eine neuen Werkstück verwendet. Dort wird sie auch zurückgesetzt. Damit sie nach dem Zurücksetzen nicht sofort wieder gesetzt wird, sollte sich das Werkstück noch unter dem Schreibkopf befinden wird mittels eines Timers \SI{6}{\second} gewartet. Das bedeutet das auch nur alle  \SI{6}{\second} ein Werkstück aus dem Lager geholt werden kann. Diese Zeit wird durch die Roboter allerdings immer weit überschritten.

\subsection{FB FB\_RFID\_SS15}
Der Funktionsbaustein\glqq FB\_RFID\_SS15\grqq wurde mit dem Starterkit zurverfügung gestellt. Der im Programm verwendete Funktionsbaustein entspricht fast kommplet dem Funktionsbaustein aus dem Starterkit weshalb er hier nicht weiter erläutert wird. Nur für das Errorhandling sind wenige Zeilen Code hinzugekommen. 

Ist der Rückgabewert "`Channel\_X\_ERROR"´ wird ein Reset durchgeführt bis der der Fehler behoben wurde. Zum Testen der Fehleranfälligkeit befindet sich ein Zähler im Errorhandling. 

\subsection{FB Werkstueck\_Tagen}
Immer wenn ein Roboter ein Werkstück aus dem Lager holen soll wird von der Fertigungsplanung in die Datenbank geschrieben um welches Werkstück es sich handelt. Dieser Funktionsbaustein fragt durchgehend die Datenbank ab und ermittelt so welches die nächste Id ist die an Station 10 auf den RFID-Tag des Werkstücks geschrieben werden muss. In Abbildung \ref{fig:FB_Werkstueck_taggen} ist das dazugehörige Zustandsdiagramm dargestellt.
\begin{figure}[h]
	    \centering
	    \includegraphics[width=0.7\linewidth]{Bilder/Zustandsdiagramme/taggen.png}
        \caption{UML-Zustandsdiagramm des FB Werkstueck\_taggen}
        \label{fig:FB_Werkstueck_taggen}
\end{figure}

\subsection{FB Timestamp\_anlegen}\label{kap:FB_Timestamp}
Der Funktionsbaustein \glqq Timestamp\_anlegen\grqq  ist zuständig für die Auswertung des RFID-Arrays, in dem die aktuell eingelesenen Werkstück id's gespeichert sind. Desweiteren Ermittelt der Funktionsbaustein alle für das anlegen des Timestamps erforderlichen Werte und schreibt mit Hilfe des Funktionsbausteins \glqq Datenbank\_Read\_Write\grqq  den Timestamp in die Datenbank. In der Abbildung \ref{fig:FB_Timestamp_anlegen} ist das Zustandsdiagramm des Funktionsbaustein dargestellt. 

\begin{figure}[h]
	    \centering
	    \includegraphics[width=1\linewidth]{Bilder/Zustandsdiagramme/Zustandsautomat_Timestamp.png}
        \caption{UML-Zustandsdiagramm des FB Timestamp\_anlegen}
        \label{fig:FB_Timestamp_anlegen}
\end{figure}

Zunächst wird im  Zustand 0 mit Hilfe der Funktion "`Su\-che\_neues\_Werk\-stueck"' nach einem neu unter den RFID-Lesekopf geschobenes Werkstück gesucht. Zu beginn wird allerdings zunächst geprüft ob an Station 10 ein Werkstück neue beschrieben wurde, dazu wird die globale Variable \glqq WRITE\_RFID\_DONE\grqq  ausgewertet und ggf. zurückgesetzt. Wurde kein Werkstück neu beschrieben wird überprüft ob sich seit dem letzten Aufruf der Funktion der Werte im RFID-Array geändert haben. Dazu wird ein Vergleich mit einer Kopie des Arrays aus dem letzten Aufruf durchgeführt. Wird eine Änderung erkannt, die sich  nicht auf 0 geändert hat, wird die Station an der die Änderung aufgetreten ist und die id des erkannten Werkstücks von der Funktion zurückgegeben. Ebenso wird die geänderte  Kopie des RFID-Arrays zurückgegeben.

Wurde ein Werkstück an Station 10 neu beschrieben, so muss zunächst mit einer Datenbank Abfrage aus der Tabelle \glqq taggen\grqq  die id des beschriebenen Werkstücks ermittelt werden. Dies geschieht im Zustand 2. Der Zustand 1 Startet die Ermittlung. Für einen Zugriff auf die Datenbank werden immer zwei Zustände benötigt, einer in dem die Abfrage gestartet wird und einer der die Abfrage ausführt und auf deren Ende wartet. 

Wurde die id des neu beschriebenen Werkstücks ermittelt, liegen genau die gleichen Informationen, nämlich die id des Werkstücks und die Station an der es erkannt wurde, wie bei einem Werkstück, dass neu unter einen Lesekopf der Stationen 2-8 geschoben wurde, vor. Für beide möglichen Fälle geht es weiter mit Zustand 3 und 4 wo die nächste frei id für den Timestamp aus der Datenbank ermittelt wird. Beim Übergang von Zustand 4 in Zustand 5 wird die aktuelle Zeit als Unix-Timestamp gespeichert. Anschließend wird im Zustand 5 und 6 der nächste Log-Status ermittelt. Der Log-Status beschreibt ob es sich um eine Anmeldung oder eine Abmeldung handelt. Zum ermitteln des Log-Status wird der letzte Eintrag für das jeweilige Werkstück in der Tabelle \glqq Timestamp\grqq  in der Datenbank ermittelt und der Status dementsprechend gesetzt. Besonders werden die Fälle an Station 10 und an Station 80 behandelt, da dort nur An- bzw. Abmeldungen möglich sind.

Nachdem nun die id des Werkstücks, die Station, die id des nächsten Timestamps, der nächste Log-Status und die aktuelle Zeit bekannt sind werden die Daten in den Zuständen 7 und 8 als Timestamp in der Datenbank abgelegt und zu Zustand 1 zurückgekehrt.

\subsection{FB Datenbank\_Read\_Write}\label{kap:FB_Datenbank_Read_Write}
Der Funktionsbaustein \glqq Datenbank\_Read\_Write\grqq  sorgt für einen geordneten Zugriff auf die Datenbank. Der eigentliche Zugriff auf die Datenbank findet in den vier Funktionsbausteinen \glqq SQL4Automation\_id\_Lesen\grqq , \glqq SQL4Automation\_nextLogStatus\-\_lesen\grqq , \glqq SQL4Automation\_RFID\_Werkstueck\_Lesen\grqq  und \glqq SQL4Automation\_Timestap\grqq \-statt. Da die Funktionsbausteine die selben globalen Variablen nutzen und mehrere Zyklen für einen Zugriff benötigen dürfen nicht mehrere der Funktionsbausteine gleichzeitig versuchen auf die Datenbank zuzugreifen. Da die Anfragen an die Datenbank von zwei verschiedenen Funktionsbausteinen aus erfolgt muss sichergestellt das die Anfragen nicht gleichzeitig erfolgt. Dafür wird dieser global instanzierte Funktionsbaustein genutzt. Wenn einer der Funktionsbausteine eine Anfrage stellen will wird geprüft ob diese Anfrage an der Reihe ist. Sollte die Anfrage nicht an der Reihe sein wartet der Funktionsbaustein bis die Anfrage an der Reihe ist und stellt jeden Zyklus erneut eine Anfrage. Der Funktionsbaustein \glqq Datenbank\_Read\_Write\grqq  prüft in dem Zustandsautomaten aus Abbildung \ref{fig:FB_Datenbank_Read_Write} ob eine Anfrage für den jeweiligen Datenbankzugriff vorliegt. Liegt keine Anfrage vor geht er zum nächsten Zustand wo auf eine Anfrage geprüft wird. Liegt einen Anfrage vor Startet er einen Datenbankzugriff und geht für die Dauer des Zugriffs in einen Zustand für die Ausführung der Anfrage. An die Anfragenden Funktionsbausteine meldet er zurück ob eine Datenbankzugriff gestartet wurde und ob der gestartete Zugriff abgeschlossen wurde.
\begin{figure}[h]
	    \centering
	    \includegraphics[width=0.8\linewidth]{Bilder/Zustandsdiagramme/Datenbank_Read_Write_State_Machine_Diagram.png}
        \caption{UML-Zustandsdiagramm des FB Datenbank\_Read\_Write}
        \label{fig:FB_Datenbank_Read_Write}
\end{figure}

\subsection{FB SQL4Automation}
Für jeden Zugriff auf die Datenbank ist ein seperater Funktionsbaustein zuständig. Aufgerufen werden die Funktionsbausteine durch den Funktionsbaustein "`Datenbank\_Read\_Write"´, der in Abschnitt \ref{kap:FB_Datenbank_Read_Write} beschrieben ist. Die Funktionsbausteine entsprechen in großen Teilen dem mit dem Staterkit zurverfügung gestellten Funktionsbaustein "`SQL4Automation\_2015"´. Angepasst wurde für die jeweiligen Funktionsbausteine der MySQL-Befehl zum Lesen oder Schreiben in der Datenbank. Desweiteren wurde bei Lesebefehlen ein zusätzlicher Rückgabewert für das Ergebnis der Abfrage und die entsprechende Konvertierung des Rückgabe-Strings der Datenbank in den richtigen Datentyp hinzugefügt. Ebenso wurde eine Rückgabevariable hinzugefügt die zurückmeldet ob die Anfrage an die Datenbank abgeschlossen ist. Wie schon in Kapitel \ref{kap:FB_Datenbank_Read_Write} beschrieben nutzen alle Funktionsbausteine, die den Zugriff auf die Datenbank durchführen die selben globalen Variablen. Da ein Zugriff mehrere Zyklen dauert ist sicherzustellen das nicht zwei Funktionsbausteine zur selben Zeit eine Anfrage an die Datenbank stellen. 

\subsection{Systemzeit Einstellen}
Da die Systemzeit nach jedem Start der Soft-SPS auf den default-Wert 1. Januar 1970 0:00:00 Uhr gesetzt wird und das manuelle einstellen der Systemzeit bei jedem Neustart der Soft-SPS Beziehungsweise des aufgespielten CODESYS-Programms ungenau und aufwendig ist wird die aktuelle Mitteleuropäische Zeit zu beginn des Programms automatisch ermittelt. Zur Ermittlung der Zeit für die Timestamps wird der Funktionsbaustein "`RTC"´ genutzt. Dieser wird zu Programmbegin im Funktionsbaustein "`Timestamp\_anlegen"´ einmal initialisiert. Für die Initialisierung wird die aktuelle Zeit im Datentype DT benötigt. Die aktuelle Zeit wird zu Programmbegin in der Main über die Datenbank ermittelt. 

Wie in Abbildung \ref{fig:ER-Diagramm_Worbenche} zu sehen gibt es in der Datenbank eine Tabelle mit dem Namen Zeit, diese Tabelle hat nur eine Zeile in der es einmal den Primärschlüssel "`id\_Zeit"` und das Feld "`aktuelle\_Zeit"´gibt. Bei Start des Programms wird in der Main zunächst der Funktionsbaustein "`SQL4Automation\_aktuelle\_Zeit"´ aufgerufen, dabei handelt es sich um den einzigen Funktionsbaustein der auf die Datenbank zugreift und nicht über den Funktionsbaustein "`Datenbank\_Read\_Write"´ aufgerufen wird. Da das ermitteln der aktuellen Zeit erst abgeschlossen sein muss bevor der restliche Programmteil startet besteht keine Gefahr eines gleichzeitigen Zugriffs.
Der Funktionsbaustein "`SQL4Automation\_aktuelle\_Zeit"´ führt die zwei in den Listings \ref{lst:MySQL_Zeit1} und \ref{lst:MySQL_Zeit2} dokumentierten SQL-Befehle aus.
Der Befehl aus Listing \ref{lst:MySQL_Zeit1} aktualisiert die Zeit im Feld "`aktuelle\_Zeit"´ in der Datenbank. Mit dem Befehl in Listing \ref{lst:MySQL_Zeit2} wird diese abgefragt.
\lstset{ 
    keywordstyle        =\bfseries\ttfamily\color{blue},
    basicstyle          =\scriptsize\ttfamily, 
    emphstyle           =\color{red},
    numbers             =left,
    xleftmargin         =15pt,
    backgroundcolor     =\color{lightgray},
    showstringspaces    =false,
    language            =SQL
    }	 

\begin{lstlisting}[caption={MySQL-Befehl: Zeit aktualisieren}
       \label{lst:MySQL_Zeit1},
       captionpos=t] 
UPDATE vpj.Zeit
SET
    aktuelle_Zeit = NOW()
WHERE
    id_ZEIT = 1;
\end{lstlisting}

\lstset{ 
    keywordstyle        =\bfseries\ttfamily\color{blue},
    basicstyle          =\scriptsize\ttfamily, 
    emphstyle           =\color{red},
    numbers             =left,
    xleftmargin         =15pt,
    backgroundcolor     =\color{lightgray},
    showstringspaces    =false,
    language            =SQL
    }	 

\begin{lstlisting}[caption={MySQL-Befehl: Zeit abfragen}
       \label{lst:MySQL_Zeit2},
       captionpos=t] 
SELECT
    aktuelle_Zeit
FROM
    vpj.Zeit
WHERE
    id_ZEIT = 1;
\end{lstlisting}

Mit Hilfe von String-Operationen wird der zurückgelieferte String so umgebaut das er mit Hilfe eines Cast vom Datentype String in den Datentype DT (Date and Time) umgewandelt werden kann. 
% !TEX root = ../VPJ.tex

\chapter{Zusammenfassung}
\label{sec:Fazit}




\listoffigures
%\listoftables


\begin{flushleft}           % kein Blocksatz, sondern Linksbündig
\bibliographystyle{alpha}
%\bibliographystyle{natdin}
%\bibliographystyle{dinat}

\bibliography{literature}
\end{flushleft}

\appendix
% !TEX root = ../VPJ.tex

\chapter{Inhalt der CD}
\begin{itemize}
    \item Dieses Dokument als PDF "`VPJ.pdf"' \\

    \item Alle konstruierten Bauteile sowohl als STL-, als auch als CATIA-Part-Datei im Ordner "`Konstruktion"' \\

    \item Den Quellcode aller Anwendungen im Ordner "`Code"'
		\begin{enumerate}
			\item config.ini
			\item config\_bandposition.py
			\item config\_farbton.py
			\item config\_referenzposition.py
			\item config\_playground.py - als Template und zum Testen von Einstellungen
		\end{enumerate}

    \item Alle Abbildungen der Arbeit im Ordner "`Abbildungen"' \\
		
    \item Ein Video "`LED\_Video.mp4"' zur Veranschaulichung der LEDs des Raspberry Pi\\
    \item Ein Video "`Raspberry Pi Test.mp4"' zur Veranschaulichung der Anwendung des Raspberry Pi auf dem Festo-Transfersystem\\
    \item Ein Video "`Pixy Test auf Anlage.mp4"' zur Veranschaulichung der Anwendung der PixyCam auf dem Festo-Transfersystem\\

    \item Marktrecherche ImageProcessingPlatform.pdf \\ "`How to choose the best embedded processing platform for on-board UAV
image processing ?"' von Dries Hulens, Jon Verbeke und Toon Goedem\\




\end{itemize}


\newpage
\clearpage
\thispagestyle{empty}
\input{style/statement_i.tex}

% appendix
\end{document}
